% ============================================================
\newpage
\section{Защитная торговля}
% ============================================================

\bidsectionstart{Балансирующая позиция}
  \bidrow{Контра}{8+}{$\sim$3 пункта «одолженных» у партнёра}
  \bidrow{Бид на 1-м ур.}{6--14}{5+ карт}
  \bidrow{\bid{1}{\nt}}{11--14}{Стоппер, сбалансированная (сниженный диапазон)}
  \bidrow{Прыжок}{12--16}{6+ карт, промежуточная рука}
\bidsectionend

\bidsectionstart{После открытия оппонентов на уровне 2}
  \bidrow{Контра}{$\sim$}{Информативная (от 10+)}
  \bidrow{Бид}{$\sim$}{Натуральный, 5+ карт}
  \bidrow{\bid{2}{\nt}}{$\sim$}{Лебенсол / натуральный (по договорённости)}
\bidsectionend

\bidsectionstart{После открытия оппонентов на уровне 3}
  \bidrow{Контра}{$\sim$}{Опционно-штрафная / информативная}
  \bidrow{\bid{3}{\nt}}{$\sim$}{К игре, стоппер}
\bidsectionend

\begin{bidnote}[Лебенсол]
  Если играете Лебенсол после контры на слабую двойку:
  \begin{itemize}[noitemsep, topsep=2pt]
    \item \bid{2}{\nt} = марионетка в \bid{3}{\cl}, далее пас или коррекция (слабая рука)
    \item Прямой бид на 3-м уровне = инвитационный (8--11)
    \item Прямое 3\nt = к игре, без стоппера
    \item Через 2\nt\ затем 3\nt = к игре, со стоппером
  \end{itemize}
  Удалите эту заметку, если не играете Лебенсол.
\end{bidnote}
